\documentclass[11pt,a4paper]{article}
\usepackage[utf8]{inputenc}
\usepackage[spanish]{babel}
\usepackage[T1]{fontenc}
\usepackage{geometry}
\geometry{margin=2.5cm}
\usepackage{titlesec}
\usepackage{enumitem}
\usepackage{longtable}
\usepackage{array}
\usepackage{xcolor}
\usepackage{hyperref}
\usepackage{fancyhdr}

\hypersetup{
    colorlinks=true,
    linkcolor=black,
    urlcolor=blue,
    citecolor=black
}

\pagestyle{fancy}
\fancyhf{}
\rhead{Reglamento Interno -- Grupo 4}
\lhead{Backend}
\rfoot{P\'agina \thepage}

\titleformat{\section}{\large\bfseries}{\thesection.}{0.5em}{}
\titleformat{\subsection}{\normalsize\bfseries}{\thesubsection.}{0.5em}{}

\title{\textbf{Reglamento Interno de Trabajo\\ Grupo 4}\\[0.3cm]
\large Asignatura: Backend}
\author{Universidad Cat\'olica de Temuco\\ Ingenier\'ia Civil Inform\'atica}
\date{21/08/2026}

\begin{document}

\maketitle
\thispagestyle{fancy}

\section{Identificaci\'on del grupo}

\begin{center}
\begin{tabular}{|l|l|}
\hline
\textbf{Integrante} & \textbf{Rol} \\
\hline
Jose Luis Sep\'ulveda & L\'ider \\
\hline
Benjam\'in Leal & Integrante \\
\hline
El\'ias Guti\'errez & Integrante \\
\hline
Lisette Delgado & Integrante \\
\hline
\end{tabular}
\end{center}

\section{Objetivo del reglamento}

El presente reglamento tiene como finalidad establecer las normas de convivencia, organizaci\'on y responsabilidad que regir\'an el trabajo del Grupo 4 durante el desarrollo de los proyectos de la asignatura Backend, con el prop\'osito de asegurar un trabajo colaborativo, equitativo y de calidad entre todos los integrantes.

\section{Del rol del l\'ider}

\begin{enumerate}[label=\alph*)]
    \item El l\'ider (Jose Luis Sep\'ulveda) ser\'a responsable de coordinar las reuniones, distribuir las tareas y velar por el cumplimiento de los plazos establecidos.
    \item El l\'ider centralizar\'a la comunicaci\'on con el profesor y ser\'a el encargado de administrar el repositorio de GitHub del grupo.
    \item El l\'ider podr\'a delegar tareas espec\'ificas, pero la responsabilidad final del cumplimiento del trabajo es compartida por todos los integrantes.
\end{enumerate}

\section{Acuerdos de trabajo}

\begin{enumerate}[label=\arabic*.]
    \item \textbf{Reuniones:} El grupo se reunir\'a, de manera presencial o v\'irtual, al menos una vez por semana durante el desarrollo de cada trabajo, en el d\'ia y horario acordado previamente por todos los integrantes.
    \item \textbf{Comunicaci\'on:} El canal oficial de comunicaci\'on del grupo ser\'a un grupo de WhatsApp, donde se informar\'an avances, dudas y cambios de horario con un m\'inimo de 24 horas de anticipaci\'on.
    \item \textbf{Cumplimiento de plazos:} Cada integrante se compromete a entregar su parte del trabajo dentro del plazo interno acordado por el grupo, el cual siempre ser\'a anterior a la fecha l\'imite de entrega del profesor, para dejar margen de revisi\'on.
    \item \textbf{Calidad del trabajo:} Todo contenido entregado debe ser elaborado por el propio integrante, revisado antes de enviarlo, y ajustado a los requerimientos indicados por el l\'ider y el profesor.
    \item \textbf{Participaci\'on equitativa:} La carga de trabajo se distribuir\'a de manera equitativa entre los cuatro integrantes, considerando las capacidades y disponibilidad de cada uno.
    \item \textbf{Presentaciones:} Todos los integrantes deben estar presentes y preparados para las instancias de presentaci\'on oral o defensa del trabajo, salvo justificaci\'on debidamente acreditada.
    \item \textbf{Respeto:} Las opiniones de todos los integrantes ser\'an escuchadas y consideradas. Las decisiones se tomar\'an preferentemente por consenso; en caso de desacuerdo, se resolver\'a por mayor\'ia simple.
\end{enumerate}

\section{M\'etodo de trabajo en GitHub}

\begin{enumerate}[label=\arabic*.]
    \item El repositorio del grupo se alojar\'a en GitHub bajo la administraci\'on del l\'ider, quien agregar\'a como colaboradores a los demás integrantes.
    \item Se utilizar\'a la rama \texttt{main} \'unicamente para las versiones finales y aprobadas por todo el grupo. El trabajo en desarrollo se realizar\'a en ramas secundarias (por ejemplo, \texttt{feature/nombre-integrante} o \texttt{feature/tarea}).
    \item Cada integrante debe subir sus avances mediante \textit{commits} con mensajes claros y descriptivos que indiquen qu\'e se modific\'o (ejemplo: \texttt{"Agrega configuraci\'on VLAN sucursal norte"}).
    \item Antes de fusionar cambios a la rama principal, se recomienda revisar el trabajo mediante un \textit{pull request}, el cual ser\'a revisado por el l\'ider u otro integrante antes de aceptarlo.
    \item El repositorio debe mantener organizada la estructura de carpetas: documentaci\'on (informe, reglamento), c\'odigo fuente del proyecto, y recursos adicionales (im\'agenes, capturas, videos).
    \item Es responsabilidad de cada integrante mantener actualizado su acceso y realizar copias de seguridad locales de su trabajo antes de subirlo.
    \item El historial de \textit{commits} de GitHub podr\'a ser utilizado como respaldo objetivo para verificar la participaci\'on real de cada integrante en caso de conflicto.
\end{enumerate}

\section{Sanciones por incumplimiento}

Con el fin de resguardar el trabajo colaborativo y la responsabilidad individual, se establecen las siguientes sanciones internas para los integrantes que incurran en incumplimientos:

\subsection{Inasistencia a reuniones}

\begin{itemize}[label=--]
    \item Primera inasistencia sin aviso previo: llamado de atenci\'on formal por parte del l\'ider, registrado por escrito en el canal oficial del grupo.
    \item Segunda inasistencia sin aviso previo: el integrante deber\'a asumir una carga adicional de trabajo, definida por el resto del grupo, como compensaci\'on.
    \item Tercera inasistencia sin aviso previo: se informar\'a la situaci\'on al profesor, dejando constancia de la reiteraci\'on del incumplimiento.
\end{itemize}

\subsection{Inasistencia a la presentaci\'on final}

\begin{itemize}[label=--]
    \item La inasistencia injustificada a la presentaci\'on o defensa oral del trabajo se considerar\'a una falta grave.
    \item En este caso, el integrante ausente podr\'a ser reportado directamente al profesor como no participante en dicha etapa del trabajo, sin perjuicio de la nota grupal obtenida por el resto del equipo.
    \item Si la inasistencia cuenta con justificaci\'on v\'alida (certificado m\'edico u otra causa acreditable) y se informa con anticipaci\'on razonable, no se aplicar\'a sanci\'on.
\end{itemize}

\subsection{No entrega de avances o incumplimiento de tareas asignadas}

\begin{itemize}[label=--]
    \item Si un integrante no entrega su parte del trabajo dentro del plazo interno acordado, sin aviso ni justificaci\'on, se le enviar\'a un recordatorio formal con un plazo final de 24 horas.
    \item Si transcurrido ese plazo el trabajo sigue sin entregarse, el resto del grupo podr\'a redistribuir la tarea y dejar constancia escrita del incumplimiento (registro en el chat grupal y en el historial de GitHub).
    \item La reiteraci\'on de incumplimientos (dos o m\'as tareas no entregadas) facultar\'a al grupo a informar la situaci\'on al profesor, solicitando que la evaluaci\'on considere la participaci\'on individual real de cada integrante.
\end{itemize}

\subsection{Trabajo de baja calidad o no realizado conforme a lo acordado}

\begin{itemize}[label=--]
    \item El trabajo entregado que no cumpla con los est\'andares m\'inimos acordados por el grupo deber\'a ser corregido por el propio integrante dentro de un plazo adicional definido por el l\'ider.
    \item De persistir la situaci\'on, se aplicar\'an las mismas medidas indicadas en la secci\'on anterior (redistribuci\'on de tareas y aviso al profesor).
\end{itemize}

\section{Resoluci\'on de conflictos}

En caso de desacuerdos o conflictos internos que no puedan resolverse mediante di\'alogo directo entre los integrantes, se seguir\'a el siguiente procedimiento:

\begin{enumerate}[label=\arabic*.]
    \item Conversaci\'on directa entre las partes involucradas, mediada por el l\'ider.
    \item Reuni\'on grupal para exponer el conflicto y buscar una soluci\'on consensuada.
    \item De no llegar a acuerdo, se solicitar\'a la intervenci\'on del profesor de la asignatura como \'ultima instancia.
\end{enumerate}

\section{Aceptaci\'on del reglamento}

Los integrantes del Grupo 4, al participar en el desarrollo de los trabajos de la asignatura Backend, declaran conocer y aceptar el presente reglamento, comprometi\'endose a cumplirlo durante todo el semestre.

\vspace{1cm}

\begin{center}
\begin{tabular}{|l|c|}
\hline
\textbf{Integrante} & \textbf{Acepta} \\
\hline
Jose Luis Sep\'ulveda (L\'ider) & \rule{2cm}{0.4pt} \\
\hline
Benjam\'in Leal & \rule{2cm}{0.4pt} \\
\hline
El\'ias Guti\'errez & \rule{2cm}{0.4pt} \\
\hline
Lisette Delgado & \rule{2cm}{0.4pt} \\
\hline
\end{tabular}
\end{center}

\end{document}